\chapter{Kubernetes jako platforma uruchomieniowa mikroserwisów}
\label{chap:kubernetes}

Kubernetes jest otwartoźródłową platformą orkiestracji kontenerów, automatyzującą ich rozmieszczanie, skalowanie
i~eksploatację na klastrze maszyn. Projekt, zainicjowany przez firmę Google w~2014 roku i~czerpiący z~jej
wieloletnich doświadczeń z~wewnętrznym systemem Borg, został przekazany fundacji Cloud Native Computing Foundation
i~stał się standardem przemysłowym w~obszarze orkiestracji \cite{burns2022}.

U~podstaw działania Kubernetes leży model deklaratywny. Użytownik zamiast wydawać platformie poleceń, opisuje
pożądany stan systemu (np. \emph{trzy repliki usługi zamówień w~wersji 1.2}). Platforma ciągle porównuje stan
faktyczny z~pożądanym. Przy różnicach niweluje rozbieżności (pętla uzgadniania, \emph{reconciliation loop}).
Z takiego modelu wynika samonaprawialność systemu. Awaria jakiegoś kontenera, węzła czy usługi powoduje rozbieżność
którą platforma próbuje usunąć.

Architektura klastra Kubernetes obejmuje płaszczyznę sterowania (\emph{control plane}), serwer API, magazyn
stanu etcd, planistę (\emph{scheduler}) oraz menedżery kontrolerów. Obejmuje również węzły robocze (\emph{worker nodes}),
na których agent \emph{kubelet} uruchamia kontenery zgodnie z~decyzjami planisty \cite{k8sdocs}. Taka architektura
pozwala usługom chmurowym wziąć odpwiedzialność za płaszczyznę sterwania. Dzięki czemu użytkownik zarządza wyłącznie
wezłami roboczymi i~uruchamianymi obciążeniami (rozdział~\ref{chap:aws}).

\section{Podstawowe obiekty Kubernetes}
\label{sec:podstawowe-obiekty-kubernetes}

Stan klastra opisywany jest za pomocą obiektów API. Poniżej omówiono obiekty wykorzystane w~zrealizowanym systemie;
ich konkretne definicje przedstawiono w~rozdziale~\ref{chap:infrastruktura}.

\textbf{Pod} jest najmniejszą jednostką wdrożeniową. Grupą jednego lub kilku kontenerów współdzielących przestrzeń
sieciową (wspólny adres IP) i~woluminy. Najczęściej w praktyce zwiera jeden kontener aplikacyjny. Pody są ulotne,
czyli w każdej chwili mogą zostać usunięte i odtworzone w innym miejscu (węźle). Dzieje się tak np podczas balansowania
gdzie Kubernetes próbuje zmieścić nowe Pody podczas wdrożenia. Wiąże się to z tym że jeśli aplikacja wymaga utrzymania stanu
to musi ją trzymać w zewnętrznym miejscu jak baza danych.

\textbf{Deployment} deklaruje pożądaną liczbę replik poda wraz z~szablonem jego specyfikacji i~strategią
aktualizacji. Za pośrednictwem obiektu ReplicaSet utrzymuje zadaną liczbę replik. Przy zmianie szablonu (np.
nowej wersji obrazu) przeprowadza aktualizację kroczącą (\emph{rolling update}). Nowe repliki uruchamiane są
stopniowo, równolegle z~wygaszaniem starych. Pozwala to wdrażać zmiany bez przerwy w~dostępności usługi.
W~zrealizowanym systemie każdy z~trzech mikroserwisów opisany jest własnym Deploymentem.

\textbf{Service} rozwiązuje problem ulotności podów, udostępniając stabilny, wirtualny punkt dostępu (nazwę DNS
i~adres IP) dla dynamicznie zmieniającego się zbioru replik, wraz z~równoważeniem obciążenia pomiędzy nimi.
Wariant ClusterIP udostępnia usługę wewnątrz klastra, NodePort --- na stałym porcie każdego węzła (wariant ten
wykorzystano do wystawienia ruchu z~zewnętrznego load balancera), a~LoadBalancer zleca utworzenie zewnętrznego
balancera dostawcy chmury.

\textbf{Namespace} dzieli klaster na logiczne przestrzenie nazw, izolując zasoby różnych zespołów lub systemów
i~stanowiąc granicę dla uprawnień oraz limitów zasobów. Komponenty zrealizowanego systemu działają w~dedykowanej
przestrzeni \code{order-portal}.

\textbf{ConfigMap i~Secret} przechowują konfigurację oddzieloną od obrazów kontenerów. ConfigMap zawiere jawną,
a~secret pofuną. Udostępniają one podom zmienne środowiskowe lub pliki. Obiekty te są infrastrukturalną realizacją 
wzorca konfiguracji zewnętrznej (sekcja~\ref{sec:wzorce-projektowe}).

\section{Zarządzanie konfiguracją i skalowaniem}
\label{sec:konfiguracja-i-skalowanie}

Specyfikacja każdego kontenera może określać zapotrzebowanie na zasoby w~postaci żądań (\emph{requests}) i~limitów
(\emph{limits}) dla procesora i~pamięci. Żądania stanowią podstawę decyzji planisty --- pod zostanie umieszczony
wyłącznie na węźle dysponującym wolnymi zasobami w~zadeklarowanej wielkości --- natomiast limity wyznaczają górną
granicę zużycia, której przekroczenie skutkuje dławieniem procesora lub zakończeniem procesu w~przypadku pamięci.
Poprawne określenie tych wartości ma szczególne znaczenie na węzłach o~małej pojemności, jakie wykorzystano
w~środowisku zrealizowanego systemu.

W specyfikacji każdego kontenera należy określić zapotrzebowanie na zasoby. Występuję one w postaci żądań (\emph{requests})
i~limitów (\emph{limits}) dla procesora i~pamięci. Na podstawie żądań planista rozmieszcza pody w węzłach
dysponującymi wolnymi zasobami w zadeklarowanej wielkości. Limity natmiast wyznaczają górną granicę zużycia, której 
przekroczenie powoduje dławieniem procesora lub zakończeniem procesu w przypadku pamięci.
Poprawne określenie tych wartości ma szczególne znaczenie na węzłach o~małej pojemności, jakie wykorzystano
w~środowisku zrealizowanego systemu.

Konfiguracja aplikacji dostarczana jest podom ze źródeł zewnętrznych. Obiektów ConfigMap i~Secret. 
W~zrealizowanym systemie poświadczenia bazy danych trafiają do podów jako zmienne środowiskowe z~obiektów Secret, 
zasilanych z~usługi AWS Secrets Manager (rozdział~\ref{chap:infrastruktura}).

Skalowanie obciążeń realizowane jest na dwóch poziomach. Poziomy pierwszy stanowi \textbf{Horizontal Pod Autoscaler}
(HPA). Kontroler, który cyklicznie porównuje zmierzone wykorzystanie zasobów (domyślnie procesora, na podstawie
danych z~komponentu \emph{metrics-server}). Z~wartością docelową i~odpowiednio zwiększa lub zmniejsza liczbę replik
Deploymentu w~zadanym przedziale \cite{k8sdocs}. Poziom drugi to skalowanie samego klastra, czyli liczby węzłów
roboczych.

\section{Trwałe składowanie danych}
\label{sec:trwale-skladowanie-danych}

System plików kontenera jest ulotny --- ginie wraz z~jego restartem --- a~pody mogą być przenoszone pomiędzy
węzłami, co czyni przechowywanie danych na dysku węzła zawodnym. Kubernetes rozwiązuje ten problem warstwą
abstrakcji nad pamięcią trwałą, rozdzielającą perspektywę administratora i~aplikacji \cite{k8sdocs}.

Zgodnie z definicją Poda, system plików kontenera jest ulotny. Pody mogą być przenoszone między węzłami.
Z tego powdu Kubernetes rozwiązuje ten problem warstwą abstrakcji nad pamięcią trwałą, 
rozdzielającą perspektywę administratora i~aplikacji \cite{k8sdocs}.

\textbf{PersistentVolume} (PV) reprezentuje fizyczny zasób pamięci --- np. wolumin blokowy dostawcy chmury ---
natomiast \textbf{PersistentVolumeClaim} (PVC) jest żądaniem aplikacji: deklaracją potrzebnej pojemności i~trybu
dostępu, niezależną od tego, skąd zasób pochodzi. \textbf{StorageClass} opisuje klasę pamięci (typ nośnika,
parametry wydajnościowe) i~umożliwia provisionowanie dynamiczne: wolumin tworzony jest automatycznie w~chwili
pojawienia się żądania, bez ręcznego przygotowywania zasobów przez administratora. Tworzeniem i~dołączaniem
woluminów zarządza sterownik zgodny ze standardem CSI (\emph{Container Storage Interface}), specyficzny dla danego
dostawcy infrastruktury.

\textbf{PersistentVolume} (PV) reprezentuje fizyczny zasób pamięci. Natomiast \textbf{PersistentVolumeClaim} (PVC) 
jest żądaniem zapotrzebowania aplikacja na taki woluminu o potrzebnej pojemności i~trybu dos†ępu.
\textbf{StorageClass} opisuje klasę pamięci (typ nośnika,parametry wydajnościowe) i~umożliwia provisionowanie dynamiczne.
Wolumin tworzony jest automatycznie w~chwili pojawienia się żądania. Tworzeniem i~dołączaniem
woluminów zarządza sterownik zgodny ze standardem CSI (\emph{Container Storage Interface}), specyficzny dla danego
dostawcy infrastruktury.

W~zrealizowanym systemie pamięci trwałej wymagają broker Kafka oraz serwer Prometheus. Wykorzystano sterownik
AWS EBS CSI z~klasą pamięci opartą na woluminach gp3, skonfigurowaną z~trybem wiązania
\emph{WaitForFirstConsumer}. Wolumin tworzony jest dopiero po zaplanowaniu poda na konkretnym węźle, co
gwarantuje jego utworzenie w~tej samej strefie dostępności, w~której działa pod. Główna baza danych systemu
pozostaje natomiast poza klastrem, jako zarządzana usługa Amazon RDS. Uzasadnienie tej decyzji przedstawiono
w~rozdziale~\ref{chap:aws}.

\section{Obserwowalność: Prometheus i Grafana}
\label{sec:obserwowalnosc-prometheus-grafana}

Obserwowalność systemu rozproszonego opiera się na trzech filarach: metrykach (liczbowych szeregach czasowych
opisujących stan systemu), logach (zapisach zdarzeń) oraz śladach rozproszonych (zapisach przebiegu pojedynczych
żądań przez kolejne usługi). W~zrealizowanym systemie zastosowano dwa pierwsze filary, z~metrykami gromadzonymi
przez Prometheus \cite{prometheusdocs}.

Prometheus działa w~modelu \emph{pull}. Aplikacje wystawiają punkty końcowe z metrykami, które cyklicznie pobiera.
Zbiór celów nie jest konfigurowany statycznie, a w~środowisku Kubernetes Prometheus odkrywa je dynamicznie poprzez API klastra.
O objęciu poda monitoringiem decydują jego adnotacje (\code{prometheus.io/scrape}, \code{prometheus.io/port}).
ozwiązanie to wykorzystano w~zrealizowanym systemie: każdy mikroserwis opatrzony jest adnotacjami wskazującymi 
port techniczny, na którym udostępnia metryki.

Metryki przyjmują postać nazwanych szeregów czasowych z~etykietami (\emph{labels}), umożliwiającymi agregację
i~filtrowanie np. liczba żądań HTTP w~podziale na ścieżkę i~kod odpowiedzi. Typowe rodzaje metryk to liczniki
(\emph{counters}), mierniki (\emph{gauges}) oraz histogramy, pozwalające szacować percentyle czasów odpowiedzi.
Zapytania do zgromadzonych danych formułowane są w~języku PromQL, który umożliwia m.in. wyznaczanie szybkości
zmian liczników czy agregacje po etykietach.

Warstwę prezentacji stanowi Grafana --- narzędzie wizualizacji, które na podstawie zapytań PromQL buduje
interaktywne kokpity (\emph{dashboards}) prezentujące stan systemu w~czasie rzeczywistym i~historycznie. W~obu
narzędziach konfigurację (źródła danych, definicje kokpitów) można dostarczać deklaratywnie, co w~zrealizowanym
systemie wykorzystano do provisionowania Grafany z~poziomu kodu infrastruktury (rozdział~\ref{chap:infrastruktura}).
Zakres metryk publikowanych przez poszczególne mikroserwisy --- w~tym metryk domenowych, takich jak liczba
utworzonych zamówień --- omówiono w~rozdziale~\ref{chap:implementacja-mikroserwisow}.

Warstwę prezentacji stanowi Grafana --- narzędzie wizualizacji, które na podstawie zapytań PromQL buduje
interaktywne kokpity (\emph{dashboards}) prezentujące stan systemu w~czasie rzeczywistym i~historycznie.
Zakres metryk publikowanych przez poszczególne mikroserwisy omówiono w~rozdziale~\ref{chap:implementacja-mikroserwisow}.

\section{Mechanizmy bezpieczeństwa w Kubernetes}
\label{sec:bezpieczenstwo-kubernetes}

Bezpieczeństwo klastra Kubernetes obejmuje kilka współdziałających warstw \cite{k8sdocs}.

\textbf{Uwierzytelnianie i~autoryzacja dostępu do API.} Każde żądanie do serwera API podlega uwierzytelnieniu,
a~następnie autoryzacji w~modelu RBAC (\emph{Role-Based Access Control}): role definiują zestawy dozwolonych
operacji na zasobach, a~powiązania ról (\emph{role bindings}) przypisują je użytkownikom lub kontom usług.
Procesy działające w~klastrze występują wobec API jako \textbf{ServiceAccount} --- konta usług o~jawnie nadawanych,
minimalnych uprawnieniach. W~zrealizowanym systemie własne konto usługi z~ograniczoną rolą posiada serwer
Prometheus, któremu do dynamicznego odkrywania celów pomiarowych wystarczą uprawnienia tylko do odczytu metadanych
podów.

Każde żądanie do serwera API pdlega uwierzytelnieniu i~autoryzacji w~mdelu RBAC (\emph{Role-Based Access Control}).
Role definiują zestawy dozwolonych operacji na zasobach. Powiązania ról (\emph{role bindings}) przypisują je użytkownikom lub kontom usług.
Procesy działające w~klastrze występują wobec API jako \textbf{ServiceAccount}. Konta usług o~jawnie nadawanych,
minimalnych uprawnieniach. W~zrealizowanym systemie własne konto usługi z~ograniczoną rolą posiada serwer
Prometheus. Psiada uprawnienia do dynamicznego odkrywania celów poprzez odczyt metadanych podów.

\textbf{Ochrona danych poufnych.} Obiekty Secret oddzielają poświadczenia od obrazów i~manifestów aplikacji,
a~dostęp do nich podlega RBAC. W~zrealizowanym systemie źródłem poświadczeń bazy danych jest AWS Secrets Manager.

\textbf{Kontekst bezpieczeństwa podów.} Specyfikacja \emph{securityContext} pozwala wymusić uruchamianie procesów
bez uprawnień administratora (\emph{nonroot}). Komplementarną praktyką jest minimalizacja samych obrazów. 
Obrazy klasy \emph{distroless} pozbawione są powłoki i~menedżera pakietów, redukują powierzchnię ataku kontenera.
Podejście to zastosowano dla wszystkich mikroserwisów systemu (rozdział~\ref{chap:implementacja-mikroserwisow}).

\textbf{Segmentacja sieci.} Domyślnie komunikacja pomiędzy podami w~klastrze jest nieograniczona. Obiekty
NetworkPolicy pozwalają zawęzić ją do jawnie dozwolonych przepływów. W~zrealizowanym systemie segmentację ruchu
zrealizowano przede wszystkim na poziomie infrastruktury chmurowej. Grup bezpieczeństwa i~topologii podsieci
(rozdział~\ref{chap:infrastruktura}).

\textbf{Integracja z~systemem uprawnień chmury.} W~środowiskach zarządzanych pody często wymagają dostępu do usług
dostawcy chmury (np. tworzenia woluminów EBS). Zamiast statycznych kluczy stosuje się mechanizmy wiążące konta
usług Kubernetes z~rolami IAM --- w~EKS są to IRSA (\emph{IAM Roles for Service Accounts}) oraz nowszy mechanizm
EKS Pod Identity, który w~zrealizowanym systemie nadaje uprawnienia sterownikowi EBS CSI. Proces otrzymuje wówczas
poświadczenia krótkoterminowe, automatycznie rotowane, o~zakresie ograniczonym do przypisanej roli.

\textbf{Integracja z~systemem uprawnień chmury.} W~środowiskach zarządzanych pody często wymagają dostępu do usług
dostawcy chmury. Wiąże się wtedy konta usług Kubernetes z rolami dostawcy. W~zrealizowanym systemie użyty został
EKS Pod Identity, który nadaje uprawnienia sterownikowi EBS CSI.

Omówione mechanizmy stanowią fundament, na którym zbudowano środowisko uruchomieniowe systemu. Sposób ich
praktycznego użycia  przedstawiono w~rozdziałach~\ref{chap:aws} i~\ref{chap:infrastruktura}.